\pagebreak

\chapter{File system}

    
    
    \begin{tabbing}
    \hspace*{2em}\=\hspace*{2em}\=\hspace*{2em}\=\hspace*{2em}\=\hspace*{2em}\=\hspace*{2em}\=\\
    \>{\bf classes/}\\
        \>\>{\bf 3pl.jar}\\
            \>\>\>3PL Java archive file\\
        \>\>{\bf threepl}\\
            \>\>\>3PL Java class files compiled from sources/threepl below.\\

    \>{\bf doc/}\\
        \>\>{\bf api}\\
            \>\>\>3PL Java API\\
        \>\>{\bf impl}\\
            \>\>\>Latex and fig sources for the 3PL implementation description\\

        \>\>{\bf overview}\\
            \>\>\>3PL overview\\

        \>\>{\bf ref}\\
            \>\>\>Latex and fig sources for the 3PL reference manual\\
    \>{\bf sources/}\\
        \>\>{\bf tools/}\\
            \>\>\>\parbox[t]{12cm}{Miscellaneous scripts and programs for configuration file generation
            or CPU to FPGA interfaces} - see \ref{ch:misctools}\\
        \>\>\>{\bf bin2c/}\\
            \>\>\>\>\parbox[t]{12cm}{Convert a binary compressed or uncompressed FPGA configuration
                bit file to a C include file containing an initialised byte
                array as C source code. By this means the configuration can be
                included in the application program and can be loaded by the CPU
                application code without the need to read an external file
                during execution.}\\

        \>\>\>{\bf bin2py/}\\
            \>\>\>\>\parbox[t]{12cm}{Convert a binary compressed or
                uncompressed FPGA configuration bit file to a Python3 source
                file containing an initialised byte array as Python3 source
                code. By this means the configuration can be included in the
                application program and can be loaded by the CPU application
                code without the need to read an external file during
                execution.}\\

        \>\>\>{\bf edif2bit/}\\
            \>\>\>\>\parbox[t]{12cm}{Generate an FPGA configuration bit file
                from the EDIF netlist file. This invokes Xilinx ISE or
                Vivado tools as appropriate.}\\
        \>\>\>{\bf info2c/}\\
                \>\>\>\>\parbox[t]{10cm}{Generate CPU/FPGA interface
                functions for C from the 3PL information file (design.info).
                This is applicable to the Zynq AXI bus interface or to
                platforms with a serial link from the CPU to the FPGA.
                This generates a C include file {\bf {\em design}\_fpga.h}}\\

        \>\>\>{\bf info2py/}\\
            \>\>\>\>\parbox[t]{12cm}{Generate CPU/FPGA interface functions for
                Python3 from the 3PL information file (design.info). This can
                be executed as a main program and will generate a Python3
                source file of CPU/FPGA interface functions that can be
                imported into the CPU application program. Alternatively this
                file can be imported into the application program in which
                case it will create the interface functions internally and
                will compile them into a module.
                
                info2py is a copy of into2py.py included so that it can be
                invoked as a main program without the file name extension, e.g.
                {\bf info2py prog} rather than python3 {\bf infotopy.py prog}.
                
                This source will import serial/cpu\_serial/cpu\_serial.py
                for FPGA platforms which have a serial interface, or
                zynq/fpga/fpga.py in the case of Zynq.}\\
                
        \>\>\>{\bf run\_ISE/}\\
            \>\>\>\>\parbox[t]{12cm}{Generate an FPGA configuration bit file
                from the EDIF netlist file using Xilinx ISE tools. This is
                a very much more primitive and direct script than edif2bit
                above.}\\

        \>\>\>{\bf run\_vivado/}\\
            \>\>\>\>\parbox[t]{12cm}{Generate an FPGA configuration bit
                file from the EDIF netlist file using Xilinx Vivado
                tools. This is a very much more primitive and direct
                script than edif2bit above.}\\

        \>\>\>{\bf serial/}\\
            \>\>\>\>\parbox[t]{12cm}{This directory contains interface code
                for FPGA platforms which use a serial interface to both
                load the FPGA configuration and to transfer data between
                the CPU and the FPGA. An example of this type of platform
                is the Gadget Factory Papilio board.}\\
                
            \>\>\>\>{\bf cpu\_serial/}\\
                    \>\>\>\>\>\parbox[t]{11cm}{Python3 low-level CPU/FPGA
                    interface communications functions.}\\
            \>\>\>\>{\bf uart\_cpu/}\\
                    \>\>\>\>\>\parbox[t]{11cm}{C low-level CPU/FPGA interface
                    communications functions.}\\
        \>\>\>{\bf zynq/}\\
            \>\>\>\>{\bf fpga/}\\
                    \>\>\>\>\>{\bf fpga.c, fpga.h}\\
                    \>\>\>\>\>\>\parbox[t]{10cm}{C low-level CPU/FPGA
                    interface communications functions.}\\                
                    \>\>\>\>\>{\bf fpga.py}\\
                    \>\>\>\>\>\>\parbox[t]{10cm}{Python3 low-level CPU/FPGA
                        interface communications. functions. This imports
                        low-level communication functions from
                        lib\_py\_ext/fpga\_read\_write.c which has been
                        added to the Python3 library by
                        lib\_py\_ext/setup.py below.}\\
            \>\>\>\>{\bf lib\_py\_ext/}\\
            \>\>\>\>\>{\bf fpga\_py\_extension.c}\\
                \>\>\>\>\>\>\parbox[t]{10cm}{C extension functions for the zynq
                        Python3 CPU/FPGA interface}\\
            \>\>\>\>\>{\bf setup.py}\\
                \>\>\>\>\>\>\parbox[t]{10cm}{A Python3 source file to add
                        fpga\_read\_write.c above as a C extension to the
                        Python3 library.}\\
        \>\>{\bf include/}\\
            \>\>\>\parbox[t]{12cm}{3PL library source files to be included in 3PL programs by
            module "\ldots" or source "\ldots" statements. These are described in
            \autorefph{ch:libraries}.}\\
            
        \>\>{\bf scripts/}\\
            \>\>\>\parbox[t]{12cm}{Scripts used by the makefile when building
            3PL.}\\
            
        \>\>{\bf threepl/}\\
            \>\>\>\parbox[t]{12cm}{Source code for the 3PL compiler and
            interpretter.}\\
            
        \>\>\>{\bf codegen/}\\
            \>\>\>\>\parbox[t]{12cm}{Java classes used to construct the
            intermediate code list.}\\
                
        \>\>\>{\bf exceptions/}\\
            \>\>\>\>\parbox[t]{12cm}{Compiler and intermediate code
            execution exceptions.}\\
                
        \>\>\>{\bf exec/}\\
            \>\>\>\>\parbox[t]{12cm}{Java classes instantiated during
                parsing and used during immediate execution.}\\
                
        \>\>\>{\bf funcs/}\\
            \>\>\>\>\parbox[t]{12cm}{3PL inbuilt functions.}\\
                
        \>\>\>{\bf mods/}\\
            \>\>\>\>\parbox[t]{12cm}{3PL inbuilt modules.}\\
                
        \>\>\>{\bf nodes/}\\
            \>\>\>\>\parbox[t]{12cm}{Nodes in the code tree produced by the
            3PL parser.
                }\\
        \>\>\>{\bf parser/}\\
            \>\>\>\>\parbox[t]{12cm}{The 3PL parser.}\\
                
        \>\>\>{\bf procs/}\\
            \>\>\>\>\parbox[t]{12cm}{3PL inbuilt procedures.}\\
                
        \>\>\>{\bf simulator/}\\
            \>\>\>\>\parbox[t]{12cm}{A target code simulator that has
                atrophied with time and neglect. It is unlikely to be still
                functional and has not been included in 3PL for quite a few
                years. It has not been deleted in case someone wants to try
                to resurrect it. This would be difficult due to lack of
                documentation!}\\
            
        \>\>\>{\bf ThreePL.java}\\
            \>\>\>\>\parbox[t]{12cm}{The 3PL main program.}\\
            
    \end{tabbing}
